\section*{Objective: Learn tens pattern and combine numbers to 100.}

\subsection*{Big Numbers 21 to 100}

\begin{vocabcard}{Tens and Big Numbers}
  \vocabpair{Twenty}{20}
  \vocabpair{Thirty}{30}
  \vocabpair{Forty}{40}
  \vocabpair{Fifty}{50}
  \vocabpair{Sixty}{60}
  \vocabpair{Seventy}{70}
\end{vocabcard}

\newpage

\begin{vocabcard}{Tens and Big Numbers (Continued)}
  \vocabpair{Eighty}{80}
  \vocabpair{Ninety}{90}
  \vocabpair{One hundred}{100}
  \vocabpair{Plus}{Más (+)}
  \vocabpair{Equals}{Igual (=)}
\end{vocabcard}

\begin{grammartip}{\faIcon{lightbulb} \ Combining Numbers}
  To make a big number, you combine the tens word with the ones word!\\
  Pattern: \textbf{Tens + ones}\\
  \emph{Example: Twenty + five = Twenty-five.}\\
  \emph{Example: Forty + two = Forty-two.}
\end{grammartip}

\begin{activity}{\faIcon{pencil-alt} \ Activity 1: Complete the Tens (Completa las decenas)}
  Write the missing tens numbers! (¡Escribe las decenas que faltan!)\\[1em]
  20, 30, \answerline[2cm], 50, \answerline[2cm], 70, \answerline[2cm], 90, \answerline[2cm].
\end{activity}

\begin{activity}{\faIcon{check} \ Activity 2: Write Big Numbers (Escribe números grandes)}
  Write these numbers in words. (Escribe estos números con letra.)\\[1em]
  1. 25 = \answerline[5cm]\\
  2. 43 = \answerline[5cm]\\
  3. 67 = \answerline[5cm]\\
  4. 81 = \answerline[5cm]\\
  5. 99 = \answerline[5cm]
\end{activity}

\begin{activity}{\faIcon{pencil-alt} \ Activity 3: Simple Math (Matemáticas simples)}
  Do the math and write the answer in words! (¡Haz la suma y escribe la respuesta con letra!)\\[1em]
  1. Thirty plus twenty = \answerline[5cm]\\
  2. Fifty plus ten = \answerline[5cm]\\
  3. Forty plus thirty = \answerline[5cm]
\end{activity}

\vspace{1em}
\begin{grammartip}{\faIcon{star} \ How did I do today?}
  Circle the stars that show how you feel about today's lesson:\\[6pt]
  \textbf{I can count by tens to 100:} \selfcheck[0] \\[4pt]
  \textbf{I can write big numbers in English:} \selfcheck[0] \\[4pt]
  \textbf{I can do simple math in English:} \selfcheck[0]
\end{grammartip}
